\chapter{How a calculation is organised}\label{ch:workflow}

A calculation runs three programs in sequence, in the same directory. All
three read the input file \file{inp}.

\begin{center}
\begin{tabular}{lll}
\toprule
program & does & writes \\
\midrule
\prog{fase1} & resonances, coefficients, band list & \file{resonanceslist}, \file{dunhamparams}, \\
 & & \file{bandlist}, \file{resonancecandidates}, \file{dumpedparams} \\
\prog{fase2a} & polyads of each starting state & \file{dir\_n*}, \file{polyadsizes} \\
\prog{fase2b} & diagonalisation, transitions & \file{hist\_*}, \file{proclog\_n*} \\
\bottomrule
\end{tabular}
\end{center}

\section{Phase 1: perturbation theory}

Phase 1 reads the harmonic frequencies, the cubic and quartic force
constants and the dipole derivatives (Chapter~\ref{ch:inputs}). It then
does three things:
\begin{itemize}
\item it decides which terms are resonant;
\item it computes all the coefficients of the transformed Hamiltonian and
  dipole;
\item it writes them to \file{dumpedparams}, a binary file read by the next
  phases.
\end{itemize}
It takes less than a minute for a molecule with about 100 modes. It also
writes three readable files that you should look at before going on:
\begin{itemize}
\item \file{resonanceslist}: the resonances found, with coupling, detuning
  and their ratio;
\item \file{dunhamparams}: the deperturbed anharmonic constants, in the
  same convention as Gaussian's X matrix (Section~\ref{sec:dunham});
\item \file{bandlist} \devonly: all the bands of the spectrum, with an intensity
  estimate and a column saying whether, and in which direction, each band
  is computed.
\end{itemize}

\section{Reviewing and adjusting (optional) \devonly}

The automatic choices of phase 1 can be changed by editing two files and
running phase 1 again. Chapter~\ref{ch:tut-review} gives an example.
\begin{itemize}
\item In \file{bandlist}, each band can be set to \texttt{both},
  \texttt{up} (absorption only), \texttt{down} (emission only) or
  \texttt{off}. Then rerun phase 1 with \opt{bandlistfile}.
\item In \file{resonancecandidates}, each candidate resonance, grouped by
  its change of quantum numbers, can be forced \texttt{on} or
  \texttt{off}. Then rerun phase 1 with \opt{resonanceoverrides}.
\end{itemize}

\section{Phase 2a: polyads}

For each starting state listed in \opt{modepopfile}, phase 2a collects the
states connected to it, and to the destinations of its bands, by the
resonant couplings into polyads. It writes their Hamiltonian matrices to
\file{dir\_n00000}, \file{dir\_n00001}, \ldots, one directory per starting
state. Its cost grows with the size of the polyads, controlled mainly by
\opt{hsizescalepar} and \opt{maxtothops}.

A polyad grows from each target state with a credit of 1. Each resonant
step spends a part of it that depends on the ratio of the coupling $W$ to
the difference $\Delta E$ of the diagonal energies of the two states; a
state is added while some credit is left \notes{§11.3}. \devonly{} With
\opt{costmodel}~$=1$ the cost of a step is $1/(\opt{costscale}\,|W/\Delta
E|+1)$, with a scale calibrated on model polyads: in one-dimensional Fermi and
Wannier--Stark ladders, the polyad then reproduces, within 1~cm$^{-1}$ and
with overlap at least 0.95, the eigenstates carrying 95\% of the target
weight. On pyrene the default \opt{costscale}~$=6$ gives polyads a
little smaller than \opt{hsizescalepar}~$=8$ (largest 3264 states against
3582, with \opt{maxtothops}~$=4$ and \opt{minhopcoupling}~$=0$). \devonly{} With the default \opt{minhopcoupling}~$=0.5$ the largest polyad of pyrene has 2363 states, and the 0~K spectrum differs from the one with \opt{minhopcoupling}~$=0$ by 0.7\% (relative L1 distance at 5~cm$^{-1}$ resolution). The cap \opt{maxtothops} stays necessary: without it the largest
polyad of pyrene has over 100\,000 states.

\section{Phase 2b: diagonalisation and spectra}

Phase 2b diagonalises the polyads (\opt{eigensolver}), or treats them
spectrally (\opt{spectral}~$=1$, for very large ones). It computes the
transitions and adds them to the histograms \file{hist\_abs\_*} and
\file{hist\_ems\_*}. These are binned in wavenumber (\opt{wnumhiststep})
and in the energy of the starting state (\opt{Ehiststep}). This is the
expensive phase. It runs with MPI and is the one to submit to a cluster.
